% answers_short.tex —— 1977 年·江苏淮阴地区初试卷 · 简短答案速查
% 由 tools/make_short_answers.py 自 answers.tex 的 \gksolhead/\gkanswer 提取，
% 并与 1977/1977_数学真题_江苏淮阴地区初试_答案要点.md 逐题比对后生成，勿手改（改源 MD 后重跑）。

\begin{gknotice}
\textbf{使用说明}：本节只给最终答案便于快速核对，完整推导见解析版。答案为编写者独立推导，\textbf{非官方标准答案}。
\end{gknotice}

\gkpart{题目}
\begin{itemize}[leftmargin=2.6em]
  \item[\textbf{1.}] (1) $8$；(2) $3-2^{\frac{13}{24}}+\sqrt{3}\left(2^{\frac{1}{6}}-2^{\frac{3}{8}}\right)\approx1.242$
  \item[\textbf{2.}] $-\dfrac{1}{2}$。
  \item[\textbf{3.}] $\dfrac{\sqrt{3}}{4}$。
  \item[\textbf{4.}] (1) $\left(2,\ -\dfrac{\pi}{6}\right)$（极角不唯一，$\left(2,\ \dfrac{11\pi}{6}\right)$ 等同价）；(2) $x^{2}+y^{2}-3x-2y=0$
  \item[\textbf{5.}] 见解析版
  \item[\textbf{6.}] 宽 $1$ 米、高 $2$ 米时通过的光线最多，最大采光面积 $2$ 平方米。
  \item[\textbf{7.}] (1) $5$；(2) $-1$
  \item[\textbf{8.}] $P(0,\ 3)$ 或 $P(0,\ -2)$。
  \item[\textbf{9.}] 见解析版
\end{itemize}
